\subsection{Primeiras missões de RVD/B}
\label{subsec_primeiras_missoes_rendezvous}

As primeiras missões de encontro e acoplamento estabeleceram os fundamentos operacionais para missões de \gls{rvdb}, em particular a necessidade de estimar estados relativos, conduzir manobras de aproximação com segurança e garantir condições geométricas adequadas para o contato final. 
Essas missões também evidenciaram, desde o início, a relevância de estratégias de guiamento e de controle capazes de lidar com incertezas, atrasos e restrições operacionais típicas de operações em órbita.

O primeiro encontro (\textit{rendezvous}) com acoplamento (\textit{docking}) entre duas espaçonaves ocorreu em $1966$, quando a tripulação do programa \emph{Gemini} realizou manualmente a aproximação e o acoplamento ao veículo não tripulado \emph{Agena} \cite{nasa1stDocking}. 
Pouco depois, em $1967$, foi demonstrada a viabilidade do acoplamento automático em órbita, com as missões soviéticas \emph{Kosmos} 186 e 188 \cite{nasaAutmaticDocking1967}. 
Esses marcos estabeleceram o contraste entre operações supervisionadas por astronautas e abordagens com maior grau de automação, além de fornecerem experiência prática sobre requisitos de alinhamento, tolerâncias de aproximação e transições seguras entre modos de operação.

Na sequência, a experiência acumulada em missões \textit{Gemini} e \textit{Apollo} ($1968$--$1972$) consolidou procedimentos de encontro, incluindo fases de aproximação progressiva, pontos de espera e decisões operacionais baseadas na avaliação contínua do estado relativo \cite{goodman2011rendezvous}. 
Esses conceitos foram posteriormente aplicados em missões como \emph{Skylab} ($1973$--$1974$) e \emph{Apollo--Soyuz}, evidenciando a maturidade das operações de encontro orbital e acoplamento como capacidade habilitadora para arquiteturas mais complexas \cite{goodman2011rendezvous}.

De forma geral, essas primeiras missões contribuíram para definir práticas que permanecem presentes em operações modernas de \gls{rvdb}, como a segmentação da aproximação em etapas, o uso de critérios de transição entre fases e a necessidade de controle robusto de atitude e de trajetória durante a fase final de aproximação. 
